\chapter{Revisão Bibliográfica}

A aplicação de aprendizado contrastivo em materiais bidimensionais (2D) ainda é um tópico pouco explorado. Na literatura, as pesquisas acerca de aprendizado contrastivo concentram-se em materiais tridimensionais (3D) ou em tarefas de classificação de imagens.
Um desafio constante na predição de propriedades físico-químicas de materiais 2D é a escassez de dados rotulados, que limita a capacidade de generalização de modelos supervisionados.  
Dessa forma, a escassez de dados rotulados para propriedades físico-químicas torna os métodos auto-supervisionados e semi-supervisionados cruciais para a análise e predição estrutural\cite{langer2022representations}. 

Esta revisão bibliográfica organiza-se em três seções principais: na Seção~\ref{sec:rb_cl}, serão discutidos os fundamentos do aprendizado contrastivo e arquiteturas recentes; na Seção~\ref{sec:rb_rs}, serão abordadas as representações estruturais de materiais via grafos e modelos multimodais; e na Seção~\ref{sec:rb_ppm2d}, serão analisadas as abordagens para predição de propriedades em materiais 2D.

\section{Aprendizado Contrastivo} \label{sec:rb_cl}

O aprendizado contrastivo (\textit{Contrastive Learning} --CL) consolidou-se como um dos paradigmas auto-supervisionados mais proeminentes no aprendizado de máquina contemporâneo. O princípio fundamental da abordagem consiste em projetar amostras em um espaço de \textit{embeddings} no qual instâncias semanticamente semelhantes (pares positivos) são aproximadas, enquanto instâncias distintas (pares negativos) são afastadas~\cite{chen2020simple}. A formulação popularizada por Chen et al.~\cite{chen2020simple} com a arquitetura SimCLR fundamenta-se na otimização da função de perda \textit{InfoNCE} (\textit{Information Noise-Contrastive Estimation}), que maximiza a informação mútua entre diferentes visões (\textit{augmentations}) geradas a partir do mesmo dado original.

Com a expansão do aprendizado de máquina para dados não euclidianos, o paradigma contrastivo foi adaptado para a aprendizagem em grafos (\textit{Graph Contrastive Learning} -- GCL). Nesse cenário, o grande desafio metodológico reside em gerar aumentações que preservem a coerência da estrutura sem introduzir ruídos significativos. Xu et al.~\cite{xu2021self} avançaram essa fronteira ao propor um arcabouço de aprendizado focado no contraste multiescala entre representações locais e globais. A abordagem de Xu et al.~\cite{xu2021self} alinha os \textit{embeddings} de sub-grafos locais com a representação global do grafo, permitindo que o modelo capture simultaneamente padrões de vizinhança imediata e a topologia geral da rede.

À medida que os métodos de GCL se diversificaram, a integridade estrutural das amostras tornou-se uma preocupação central. Chen et al.~\cite{chen2024topogcl} desenvolveram o \textit{TopoGCL} (\textit{Topological Graph Contrastive Learning}), incorporando invariantes topológicos e análise de dados ao processo contrastivo. Chen et al.~\cite{chen2024topogcl} demonstraram que guiar o agrupamento no espaço latente preservando a topologia intrínseca do grafo impede o colapso de representação e melhora substancialmente a generalização do modelo em tarefas \textit{downstream}.

Além da vertente puramente auto-supervisionada, a literatura consolidou variações de aprendizado contrastivo supervisionado. Nessa abordagem, o espaço latente é organizado utilizando rótulos conhecidos para orientar a formação dos pares positivos e negativos. Enquanto o contraste auto-supervisionado busca uma representação genérica baseada apenas em invariâncias estruturais, a variante supervisionada força a clusterização de amostras que compartilham propriedades semelhantes, sendo especialmente útil em cenários em que o objetivo final exige uma regressão de alta precisão.

A transição de um paradigma supervisionado rígido, cujas limitações de generalização diante da escassez de dados rotulados foram evidenciadas no trabalho de \textit{benchmark} MoleculeNet de Wu et al.~\cite{wu2018moleculenet}, para o uso de \textit{encoders} contrastivos viabilizou o surgimento dos \textit{Foundation Models}. Conforme sintetizado no \textit{survey} de Van et al.~\cite{van2026survey}, o pré-treinamento contrastivo em grandes volumes de dados não rotulados permite que a rede \textit{encoder} aprenda descritores universais. Essa representação pré-treinada pode, subsequentemente, ser transferida e refinada (\textit{fine-tuning}) para tarefas específicas com conjuntos de dados reduzidos.

No paradigma de redes siamesas aplicadas ao aprendizado contrastivo, duas cópias de uma mesma rede codificadora $f_\theta(\cdot)$ com pesos compartilhados projetam um par de entradas $(\mathbf{x}_i, \mathbf{x}_j)$ para vetores latentes $\mathbf{z}_i = f_\theta(\mathbf{x}_i)$ e $\mathbf{z}_j = f_\theta(\mathbf{x}_j)$. A similaridade entre essas representações é frequentemente mensurada pela similaridade do cosseno:

\begin{equation}
    \text{sim}(\mathbf{z}_i, \mathbf{z}_j) = \frac{\mathbf{z}_i^\top \mathbf{z}_j}{\|\mathbf{z}_i\|_2 \, \|\mathbf{z}_j\|_2}
\end{equation}

Com base nessa métrica (ou na distância do cosseno $d_{\cos} = 1 - \text{sim}(\mathbf{z}_i, \mathbf{z}_j)$), funções de perda contrastivas — como a \textit{Contrastive Loss} baseada em margem ou a perda NT-Xent (\textit{Normalized Temperature-scaled Cross Entropy}) — ajustam os parâmetros $\theta$ de modo a aproximar pares positivos no espaço latente e afastar pares negativos.


\section{Representações Estruturais de Materiais} \label{sec:rb_rs}

A modelagem de propriedades físico-químicas via aprendizado de máquina requer a conversão de arranjos atômicos e geometrias cristalinas em descritores computacionais capazes de alimentar redes neurais. Historicamente, essa representação dependia de descritores estruturais desenhados manualmente (\textit{engineered fingerprints}), como as matrizes de Coulomb, representações de tensores de muitos corpos (\textit{Many-Body Tensor Representation} - MBTR) e o método SOAP (\textit{Smooth Overlap of Atomic Positions})~\cite{bartok2013representing}. O algoritmo SOAP, por exemplo, mapeia o ambiente atômico local através de distribuições contínuas de densidade gaussiana expansíveis em harmônicos esféricos, garantindo invariância por rotação, translação e permutação de átomos idênticos~\cite{bartok2013representing}. Embora eficientes quando acoplados a modelos tradicionais como \textit{Kernel Ridge Regression} (KRR), descritores artesanais apresentam alta dimensão computacional e rigidez, falhando em capturar interações de longo alcance e simetrias globais complexas de forma adaptativa.

% Para alimentar modelos de aprendizado de máquina com estruturas cristalinas periódicas, é necessário garantir invariância à translação, rotação e permutação de átomos idênticos. A Representação Tensorial de Muitos Corpos (\textit{Many-Body Tensor Representation} -- MBTR) codifica a estrutura atômica por meio de distribuições contínuas ponderadas de termos geométricos de diferentes ordens: o termo de um corpo ($K_1$) descreve a composição química (números atômicos presentes na célula unitária), o termo de dois corpos ($K_2$) captura a distribuição das distâncias inversas entre pares de átomos, e o termo de três corpos ($K_3$) representa os ângulos entre trios atômicos. Como vetores MBTR costumam apresentar alta dimensionalidade e esparsidade, técnicas de seleção de atributos, como filtros de variância, são empregadas para descartar dimensões pouco informativas antes do treinamento.

A superação das limitações dos descritores estáticos ocorreu com a introdução das redes neurais baseadas em grafos (\textit{Graph Neural Networks} - GNNs), que consolidaram a aprendizagem profunda de representações estruturais como o estado da arte para a predição de propriedades de materiais~\cite{xie2018crystal}. Sob essa perspectiva, uma estrutura cristalina é mapeada naturalmente como um grafo tridimensional $G = (V, E)$, no qual os nós ($V$) representam os átomos (codificados por vetores de características periódicas) e as arestas ($E$) representam as ligações químicas ou distâncias interatômicas relativas. Arquiteturas pioneiras, como a CGCNN (\textit{Crystal Graph Convolutional Neural Networks}) de Xie e Grossman~\cite{xie2018crystal}, demonstraram que passar mensagens convolucionais entre nós vizinhos permite aprender representações ponta a ponta (\textit{end-to-end}) diretamente a partir das coordenadas e das condições periódicas de fronteira, superando os fingerprints tradicionais em termos de acurácia e generalização.

Atualmente, o estado da arte na aprendizagem em grafos orientados à física é dominado pelo paradigma do Aprendizado Profundo Geométrico (\textit{Geometric Deep Learning}) e pelas redes equivariantemente ajustadas ao grupo de isometria tridimensional $E(3)$ ou $SE(3)$. Modelos como SchNet~\cite{schutt2018schnet}, DimeNet~\cite{klicpera2020directional} e, notadamente, a família Equiformer~\cite{liao2022equiformer} incorporam invariância e equivariância translacional e rotacional diretamente nas funções de atualização dos nós. O Equiformer~\cite{liao2022equiformer} aproveita a arquitetura de Transformers em conjunto com representações de harmônicos esféricos de alta ordem, permitindo capturar interações tensorias e direcionais complexas de múltiplos corpos. Essa formulação equivariante garante que transformações rígidas na célula unitária não alterem as predições físicas, estabelecendo o desempenho de ponta em benchmarks de propriedades termodinâmicas, eletrônicas e mecânicas.

No âmbito da aprendizagem auto-supervisionada, a eficiência das representações em grafos depende da definição de aumentações que preservem a coerência física do sistema cristalino. Koker et al.~\cite{koker2022graph} desenvolveram o arcabouço \textit{CrystalCLR}, aplicando aprendizado contrastivo sobre redes CGCNN treinadas com estruturas 3D provenientes do Materials Project. O trabalho de Koker et al.~\cite{koker2022graph} demonstrou que transformações domain-specific aplicadas à malha cristalina, tais como perturbações atômicas controladas e deformações elásticas na célula unitária, permitem que o encoder de grafos aprenda descritores estruturais latentes que superam significativamente os descritores manuais em tarefas downstream de regressão de propriedades, em específico, para temperatura de fusão, condutividade térmica e módulo volumétrico.
Pinheiro et al.~\cite{Pinheiro2022} também exploraram o aprendizado contrastivo em grafos para a predição de propriedades moleculares na base de dados QM9. 
Ao utilizar SMILES e grafos moleculares como representação de entrada, foi possível aplicar transformações nas representações e utilizar diferentes \textit{encoders} para mapear representações do mesmo composto para regiões próximas no espaço latente.
Ambos trabalhos servirão de alicerce para o desenvolvimento do modelo proposto neste trabalho, que visa adaptar o arcabouço de aprendizado contrastivo para materiais bidimensionais, considerando as particularidades geométricas e topológicas das monocamadas.



Por fim, a evolução mais recente na representação de materiais busca integrar múltiplos canais de informação estrutural por meio de encoders multimodais e Foundation Models~\cite{van2026survey}. Um exemplo proeminente é o modelo MultiMat (Moro et al., 2025)~\cite{moro2025multimodal}, que adapta o conceito de alinhamento latente da arquitetura CLIP ao domínio dos materiais. O MultiMat projeta diferentes modalidades do mesmo material como grafos atômicos via GNNs, densidades eletrônicas volumétricas via 3D-CNNs e propriedades textuais via Transformers — em um espaço de embeddings compartilhado. A convergência entre redes equivariantes de alto desempenho (como o Equiformer), transformações estruturais fisicamente orientadas (CrystalCLR) e espaços latentes multimodais (MultiMat) fundamenta as fronteiras atuais para a representação eficiente de materiais.

\section{Predição de Propriedades de Materiais 2D} \label{sec:rb_ppm2d}

Os materiais bidimensionais (2D) possuem propriedades físico-químicas, eletrônicas e mecânicas singulares resultantes de seus efeitos de confinamento quântico e reduzida espessura atômica, tornando-os promissores para aplicações em microeletrônica, catálise, fotovoltaica e armazenamento de energia. Contudo, a caracterização dessas propriedades via cálculos de Primeiros Princípios baseados na Teoria do Funcional da Densidade (DFT) apresenta alto custo computacional, especialmente na estimativa precisa do \textit{bandgap} ($E_g$) e da energia de formação ($\Delta H$), o que limita a varredura em larga escala do espaço composicional e estrutural.

Embora métodos de aprendizado supervisionado em grafos tenham alcançado alta acurácia em materiais tridimensionais (3D), a sua aplicação direta em sistemas 2D enfrenta o gargalo da escassez de dados rotulados. A limitação de amostras sintetizadas experimentalmente ou calculadas computacionalmente via métodos de alto nível (como funcionais híbridos HSE06 ou aproximação GW) restringe a capacidade de generalização de redes neurais profundas. Esse cenário justifica a adoção de paradigmas auto-supervisionados e contrastivos, os quais utilizam estruturas não rotuladas para construir um espaço latente bem condicionado, minimizando o \textit{overfitting} durante a fase posterior de ajuste fino (\textit{fine-tuning}).

A escolha dos conjuntos de dados de treinamento é determinante para a fidelidade dos modelos preditivos. Enquanto bases de dados 3D massivas, como o \textit{Materials Project}~\cite{jain2013commentary}, contam com centenas de milhares de estruturas, a qualidade e o tipo dos dados variam consideravelmente. Como observado por Moro et al.~\cite{moro2025multimodal}, o \textit{Materials Project} apresenta limitações na precisão do \textit{bandgap} devido ao uso extensivo do funcional PBE (que subestima sistematicamente a lacuna de faixa proibida), levando à necessidade de recorrer a bases especializadas como o \textit{SNUMAT} ou a dados refinados via HSE06/GW.

No âmbito dos materiais monocamada, destacam-se a \textit{Computational 2D Materials Database} (C2DB)~\cite{haastrup2018computational, gantenbein2021c2db} e o repositório \textit{JARVIS-2D}~\cite{choudhary2020joint}. A C2DB compila milhares de monocamadas 2D sistematicamente caracterizadas sob protocolos DFT padronizados, contendo tanto a energia de formação quanto valores de \textit{bandgap} obtidos por PBE e HSE06. Estudos recentes apontam que a combinação de múltiplos conjuntos de dados, como o pré-treinamento em bases 3D (ex: \textit{Materials Project}) seguido do aprendizado por transferência (\textit{transfer learning}) para bases 2D (C2DB e MC2D), melhora expressivamente a predição em regime de dados escassos, superando modelos treinados exclusivamente em uma única base isolada.

Atualmente, o estado da arte na predição de propriedades em materiais cristalinos e 2D é liderado por arquiteturas de grafos que incorporam informações de linha (Line Graph Neural Networks) e invariâncias geométricas. Redes como o ALIGNN (\textit{Atomistic Line Graph Neural Network})~\cite{choudhary2021atomistic} e variações equivariantes baseadas em Transformers (como Equiformer~\cite{liao2022equiformer}) alcançaram Erros Médios Absolutos (MAE) inferiores a $0.05$ ev/átomo para a energia de formação e da ordem de $0.15-0.25$ eV para o \textit{bandgap} na base C2DB e no repositório JARVIS.A integração dessas redes de alto desempenho (encoders ALIGNN ou Equiformer) com estratégias de aprendizado contrastivo pré-treinado no C2DB estabelece o ecossistema ideal para otimizar conjuntamente a predição da energia de formação e do \textit{bandgap} a partir de uma representação latente unificada. 

% No quesito de aprendizado contrastivo aplicado a predição de propriedades, o trabalho de Koker et al.~\cite{koker2022graph} demonstrou que a aplicação de aumentações (\textit{augmentations}) físicas específicas em estruturas cristalinas 3D permite que o modelo aprenda representações latentes invariantes que superam significativamente os descritores manuais em tarefas downstream de regressão de propriedades, Para . A adaptação dessas técnicas para materiais 2D, considerando as particularidades geométricas e topológicas das monocamadas, constitui um desafio metodológico relevante e uma oportunidade de pesquisa promissora.

